\section{Matematik mayatnik}

\begin{enumerate}
    \item Tebranish davri 0,5 s bo‘lgan matematik mayatnikning uzunligini toping (cm). $g = 10 \text{ m/s}^2, \pi^2 \approx 10$.
    \begin{itemize}
        \item[A)] 4,15.
        \item[B)] 25.
        \item[C)] 6,25.
        \item[D)] 12,5.
        \item[E)] 8,1.
    \end{itemize}

    \item Matematik mayatnikning uzunligi 2,5 m, unga osilgan sharchaning massasi 100 g. Tebranish davri qanday (s)?
    \begin{itemize}
        \item[A)] 3.14.
        \item[B)] 1,57.
        \item[C)] 2,5.
        \item[D)] 3.
    \end{itemize}

    \item Matematik mayatnikning tebranish davri nimalarga bog‘liq?
    1) mayatnik massasiga; 2) tebranish amplitudasiga; 3) erkin tushish tezlanishiga; 4) mayatnik uzunligiga.
    \begin{itemize}
        \item[A)] 1, 4.
        \item[B)] 2, 4.
        \item[C)] 2, 3.
        \item[D)] 3, 4.
        
    \end{itemize}

    \item Agar tebranishlar davri $T$ bo‘lgan matematik mayatnikning uzunligi $n$ marta oshirilsa, tebranishlar davri qanday bo‘ladi?
    \begin{itemize}
        \item[A)] $nT$
        \item[B)] $2n T$
        \item[C)] $T\sqrt{n}$
        \item[D)] $\sqrt{n}T$
        
    \end{itemize}

    \item Ipining uzunligi $l$ bo‘lgan matematik mayatnikning tebranish davri $T$ ga teng. Agar ipning $1/2$ qismini kesib tashlab, qolgan qismiga yana o‘sha yuk osilsa, tebranish davri qanday bo‘ladi?
    \begin{itemize}
        \item[A)] $T/4$.
        \item[B)] $\sqrt{2}T$.
        \item[C)] $T/\sqrt{2}$.
        \item[D)] $T/2$.
        
    \end{itemize}

    \item Matematik mayatnikning uzunligi va tebranishlari amplitudasi 16 marta oshirilsa, uning tebranishlari chastotasi qanday o‘zgaradi?
    \begin{itemize}
        \item[A)] 4 marta ortadi.
        \item[B)] 2 marta ortadi.
        \item[C)] 2 marta kamayadi.
        \item[D)] 4 marta ortadi.
        
    \end{itemize}

    \item Lift 6 m/s$^2$ tezlanish bilan pastga tushmoqda. Liftga osilgan 1 m uzunlikdagi matematik mayatnikning tebranish davrini toping (s). $g=10 \text{ m/s}^2$.
    \begin{itemize}
        \item[A)] 6,28.
        \item[B)] 2,1.
        \item[C)] 1.
        
        \item[D)] 1,57.
    \end{itemize}

    \item Matematik mayatnik uzunligi 8 cm ortganda, tebranishlar davri 0,1 s oshdi. Tebranishlarning dastlabki davri qanday (s)? ($\pi^2 \approx g$).
    \begin{itemize}
        \item[A)] 1,4.
        \item[B)] 1,6.
        \item[C)] 1,7.
        \item[D)] 2.
        
    \end{itemize}

    \item Yuqoriga 4 m/s$^2$ tezlanish bilan tekis tezlanuvchan ko‘tarilayotgan 126 cm uzunlikdagi matematik mayatnikning siklik chastotasini toping (rad/s).
    \begin{itemize}
        \item[A)] 20/9.
        \item[B)] 12/5.
        \item[C)] 10/3.
        \item[D)] 14/3.
        
    \end{itemize}

    \item Massasi 50 g bo‘lgan po‘lat sharchali mayatnikning tebranish davri 2 s. Sharcha ostiga magnit joylashtirilganda, mayatnikning tebranish davri 1 s gacha kamaydi. Sharchaning magnitga tortilish kuchi qanday (N)?
    \begin{itemize}
        \item[A)] 0,5.
        \item[B)] 1.
        \item[C)] 5.
        \item[D)] 3.
        
    \end{itemize}

    \item Birinchi mayatnikning tebranish davri 8 s, ikkinchisiniki 6 s. Ularning uzunliklari yig‘indisiga teng uzunlikdagi matematik mayatnikning tebranish davri necha sekund?
    \begin{itemize}
        \item[A)] 8.
        \item[B)] 12.
        \item[C)] 14.
        \item[D)] 5.
        
    \end{itemize}

    \item Bir vaqtda ikkita matematik mayatnikning biri 20 marta, ikkinchisi 30 marta tebrandi. Ularning uzunliklari farqi 10 cm bo‘lsa, uzunliklari yig‘indisi qanday bo‘ladi (cm)?
    \begin{itemize}
        \item[A)] 30.
        \item[B)] 20.
        \item[C)] 26.
        \item[D)] 32.
        
    \end{itemize}

    \item $60^\circ$ burchakka og‘dirilgan, massasi 50 g bo‘lgan mayatnikni muvozanat vaziyatiga qaytaruvchi kuch qanday (N)?
    \begin{itemize}
        \item[A)] $0,25\sqrt{3}$.
        \item[B)] $0,5\sqrt{3}$.
        \item[C)] 2.
        \item[D)] 3.
        
    \end{itemize}
\end{enumerate}




\begin{center}
\begin{tabular}{|c|c|c|c|c|}
\hline
1 & 2 & 3 & 4 & 5 \\
\hline
D & A & B & D & C \\
\hline
6 & 7 & 8 & 9 & 10 \\
\hline
A & D & A & C & A \\
\hline
11 & 12 & 13 & & \\
\hline
D & A & A & & \\
\hline
\end{tabular}
\end{center}